% Background chapter: definitions, examples, and setup.

\section{Background}\label{sec:background}

This chapter sets up the objects the paper reasons about. It begins with
graphs and features (Section~\ref{sec:graphs}), then defines colour refinement
(Section~\ref{sec:wl}) and message-passing graph neural networks
(Section~\ref{sec:mpnn}), and closes with the two function classes the
dichotomy of Section~\ref{sec:dichotomy} compares (Section~\ref{sec:classes}).

\subsection{Graphs and features}\label{sec:graphs}

A \emph{graph} is a pair $\gph = (\vertices, \edges)$ in which
$\vertices$ is a finite set of \emph{vertices} and $\edges \subseteq
\binom{\vertices}{2}$ is a set of \emph{edges}. Every graph in this
paper is finite, undirected, and simple: no loops from a vertex to
itself, no multi-edges, and every edge is an unordered pair of distinct
vertices. Write $|\vertices|$ for the number of vertices and $|\edges|$
for the number of edges. The \emph{degree} of a vertex $v$ is the
number of edges containing it: $\deg_\gph(v) = |\{u \in \vertices :
\{u, v\} \in \edges\}|$. The \emph{neighbourhood} of $v$ is
\[
\neigh{v} = \{u \in \vertices : \{u, v\} \in \edges\}.
\]

Two graphs $\gph = (\vertices, \edges)$ and $\gph' = (\vertices',
\edges')$ are \emph{isomorphic}, written $\gph \cong \gph'$, if there
is a bijection $\pi : \vertices \to \vertices'$ that maps edges to
edges: $\{u, v\} \in \edges$ if and only if $\{\pi(u), \pi(v)\} \in
\edges'$. Isomorphism is the natural equivalence for graphs
considered as combinatorial objects: two isomorphic graphs are the
same graph up to renaming of vertices.

A \emph{feature} on a graph $\gph$ is a map $x : \vertices \to
\mathcal{X}$ from vertices to a set $\mathcal{X}$ called the
\emph{feature space}. The paper works with $\mathcal{X} \subset \R^d$
for some $d$; the specific $d$ varies by context. A \emph{featured
graph} is a pair $(\gph, x)$. When the feature space $\mathcal{X}$ has
a single element (or is left unspecified), the graph is
\emph{unfeatured}.

Two featured graphs $(\gph, x)$ and $(\gph', x')$ are
\emph{isomorphic}, written $(\gph, x) \cong (\gph', x')$, if there is
a bijection $\pi : \vertices \to \vertices'$ that maps edges to edges
and preserves features: $x'(\pi(v)) = x(v)$ for every $v \in
\vertices$. The natural equivalence on featured graphs is featured
isomorphism.

\begin{example}\label{ex:graphs}
Let $\vertices = \{1, 2, 3, 4, 5, 6\}$. The graph $A$ with edges
$\{\{1,2\}, \{1,3\}, \{2,3\}, \{4,5\}, \{4,6\}, \{5,6\}\}$ is the
disjoint union of two triangles, written $A = C_3 \sqcup C_3$. Every
vertex has degree $2$. The graph $B$ with edges $\{\{1,2\}, \{2,3\},
\{3,4\}, \{4,5\}, \{5,6\}, \{1,6\}\}$ is the six-cycle $C_6$. Every
vertex of $B$ also has degree $2$. The two graphs are not isomorphic:
$A$ has two connected components while $B$ has one. Both graphs will
serve as running examples throughout the paper.
\end{example}

\subsection{Colour refinement (1-WL)}\label{sec:wl}

Colour refinement is an algorithm that assigns colours to the vertices
of a featured graph so that vertices in the same colour class have
locally indistinguishable neighbourhoods. It is a subroutine of the
Weisfeiler-Leman graph isomorphism test in its $1$-dimensional form,
and we use the abbreviation $1$-WL for the refinement. The relevance
to graph neural networks is that a wide class of message-passing
architectures produces vertex representations no more expressive than
the colours refinement assigns; the dichotomy in
Section~\ref{sec:dichotomy} makes this precise.

\begin{definition}[Colour refinement]\label{def:wl}
Given a featured graph $(\gph, x)$, colour refinement produces
per-vertex colours $c_v^{(t)} \in \Sigma$ for $v \in \vertices$ and
$t \in \{0, 1, 2, \dots\}$ by
\[
c_v^{(0)} = \tau(x_v),
\qquad
c_v^{(t)} = \tau\!\left(c_v^{(t-1)},\ \{\!\{c_u^{(t-1)} : u \in
\neigh{v}\}\!\}\right),
\]
where $\Sigma$ is a countable alphabet, $\{\!\{\cdot\}\!\}$ denotes a
multiset, and $\tau$ is any injective hash from its input into
$\Sigma$.
\end{definition}

The multiset $\{\!\{c_u^{(t-1)} : u \in \neigh{v}\}\!\}$ records the
colours of the neighbours of $v$, with their multiplicities but not
their identities. Two vertices with the same colour at round $t-1$
and neighbourhoods carrying the same multiset of colours receive the
same colour at round $t$. The injective hash $\tau$ separates any two
distinct inputs into distinct output colours.

The refinement stabilises: there is a round $t^\ast$ after which
$c_v^{(t)} = c_v^{(t^\ast)}$ for every $v$ and every $t \ge t^\ast$.
The reason is that the equivalence relation "$c_v^{(t)} = c_u^{(t)}$"
on vertices refines from round to round; a graph on $|\vertices|$
vertices admits at most $|\vertices|$ distinct equivalence classes, so
the refinement stops at $t^\ast \le |\vertices|$. Write $c_v =
c_v^{(t^\ast)}$ for the stable colour of $v$, and
\[
\mathrm{hist}(\gph, x) = \{\!\{c_v : v \in \vertices\}\!\}
\]
for the graph-level \emph{fingerprint}. Two featured graphs $(\gph,
x)$ and $(\gph', x')$ are \emph{$1$-WL equivalent}, written $(\gph,
x) \sim_{\wl} (\gph', x')$, if $\mathrm{hist}(\gph, x) = \mathrm{hist}(\gph',
x')$.

\begin{example}[Refinement on the six-cycle]\label{ex:wl-c6}
Take the six-cycle $B = C_6$ from Example~\ref{ex:graphs} with the constant
feature $x_v \equiv \bot$. Colour refinement runs as follows.
\begin{itemize}
\item Round $0$: $c_v^{(0)} = \tau(\bot) = \alpha$ for every $v$.
\item Round $1$: every vertex has two neighbours, both coloured
$\alpha$, so $c_v^{(1)} = \tau(\alpha, \{\!\{\alpha, \alpha\}\!\}) =
\beta$ for every $v$.
\item Round $2$: every vertex has two neighbours coloured $\beta$, so
$c_v^{(2)} = \tau(\beta, \{\!\{\beta, \beta\}\!\}) = \gamma$ for every
$v$.
\end{itemize}
All colour classes at every round have size six. Refinement stabilises
after round one (the assignment of a single colour to all vertices is
the coarsest possible partition). The fingerprint is
$\mathrm{hist}(B, \bot) = \{\!\{\beta^6\}\!\}$ (six copies of $\beta$).

The same computation on $A = C_3 \sqcup C_3$ (with the same constant
feature) also produces a single colour with multiplicity six at
round one, and $\mathrm{hist}(A, \bot) = \mathrm{hist}(B, \bot)$.
Hence $A \sim_{\wl} B$: colour refinement cannot distinguish the
disjoint-triangles graph from the six-cycle.
\end{example}

\begin{example}[Refinement with distinct features]\label{ex:wl-feat}
Take the same graph $B = C_6$ with vertices in order around the cycle
$1, 2, 3, 4, 5, 6$. Assign features $x_1 = x_3 = x_5 = +1$ and
$x_2 = x_4 = x_6 = -1$. Then
\begin{itemize}
\item Round $0$: $c_v^{(0)}$ is $\alpha_+$ on odd vertices and
$\alpha_-$ on even vertices.
\item Round $1$: every odd vertex has two neighbours coloured
$\alpha_-$, so $c_v^{(1)} = \tau(\alpha_+, \{\!\{\alpha_-,
\alpha_-\}\!\}) = \beta_+$. Every even vertex has two neighbours
coloured $\alpha_+$, so $c_v^{(1)} = \tau(\alpha_-, \{\!\{\alpha_+,
\alpha_+\}\!\}) = \beta_-$.
\item Round $2$: neighbour multisets are $\{\!\{\beta_-,
\beta_-\}\!\}$ at odd vertices and $\{\!\{\beta_+, \beta_+\}\!\}$ at
even ones; refinement stabilises at $\{\!\{\beta_+^3,\ \beta_-^3\}\!\}$
after round one.
\end{itemize}
Features preserve their role at every round: the alternating feature
pattern is visible to the refinement and propagates upward.
\end{example}

\begin{example}[Refinement on a small tree]\label{ex:wl-tree}
Consider the path graph $P_4$ with four vertices $\{1, 2, 3, 4\}$
and three edges $\{1, 2\}, \{2, 3\}, \{3, 4\}$. Every vertex has an
unfeatured constant $\bot$. The initial colouring is $c_v^{(0)} =
\alpha$ everywhere.
\begin{itemize}
\item Round $1$: vertex $1$ has one neighbour (vertex $2$), so
$c_1^{(1)} = \tau(\alpha, \{\!\{\alpha\}\!\}) = \beta_1$. Vertex $4$
also has one neighbour, so $c_4^{(1)} = \beta_1$. Vertex $2$ has two
neighbours (vertices $1$ and $3$), so $c_2^{(1)} = \tau(\alpha,
\{\!\{\alpha, \alpha\}\!\}) = \beta_2$; similarly $c_3^{(1)} =
\beta_2$.
\item Round $2$: vertex $1$'s only neighbour has colour $\beta_2$,
so $c_1^{(2)} = \tau(\beta_1, \{\!\{\beta_2\}\!\}) = \gamma_1$; same
for vertex $4$. Vertex $2$'s neighbours have colours $\beta_1$ and
$\beta_2$, so $c_2^{(2)} = \tau(\beta_2, \{\!\{\beta_1, \beta_2\}\!\})
= \gamma_2$; same for vertex $3$.
\item Round $3$: applied to $\gamma_1$ and $\gamma_2$ colours; the
refinement partition does not further split (still $\{\{1, 4\},
\{2, 3\}\}$).
\end{itemize}
$P_4$'s $1$-WL fingerprint is $\{\!\{\gamma_1, \gamma_1, \gamma_2,
\gamma_2\}\!\}$. Note that the endpoints (degree $1$) and the
interior vertices (degree $2$) receive distinct stable colours: the
refinement distinguishes them from round $1$ onward.
\end{example}

\begin{example}[Refinement on a star]\label{ex:wl-star}
The star graph $K_{1,4}$ has a centre vertex $c$ and four leaves
$l_1, l_2, l_3, l_4$; edges are $\{c, l_i\}$ for $i \in \{1, 2, 3,
4\}$. Under constant feature $\bot$:
\begin{itemize}
\item Round $0$: every vertex has colour $\alpha$.
\item Round $1$: the centre has four neighbours all coloured
$\alpha$, so $c_c^{(1)} = \tau(\alpha, \{\!\{\alpha, \alpha, \alpha,
\alpha\}\!\}) = \gamma_c$. Each leaf has one neighbour with colour
$\alpha$, so $c_{l_i}^{(1)} = \tau(\alpha, \{\!\{\alpha\}\!\}) =
\gamma_l$.
\item Round $2$: the centre's neighbours all have colour $\gamma_l$,
so $c_c^{(2)}$ is a new colour distinct from any leaf's colour.
Each leaf's only neighbour has colour $\gamma_c$, so all leaves
share a common new colour distinct from the centre's.
\end{itemize}
The refinement partition stabilises at $\{\{c\}, \{l_1, l_2, l_3,
l_4\}\}$: the centre is separated from the leaves. Fingerprint:
$\{\!\{\gamma_c^1, \gamma_l^4\}\!\}$ (one centre colour, four leaf
colours).
\end{example}

\begin{example}[1-WL distinguishes $P_4$ from $K_{1,3}$]\label{ex:wl-distinguish}
Consider $P_4$ from Example~\ref{ex:wl-tree} and $K_{1,3}$ (a star with
three leaves, on four vertices, three edges). Both graphs have four
vertices, three edges, and are trees. Under constant feature $\bot$:
$P_4$'s degree sequence is $(1, 2, 2, 1)$; $K_{1,3}$'s is $(3, 1, 1,
1)$. Round $1$ of $1$-WL separates by degree, so $P_4$'s fingerprint
at round $1$ is $\{\!\{\beta_1, \beta_2, \beta_2, \beta_1\}\!\}$ and
$K_{1,3}$'s is $\{\!\{\beta_3, \beta_1, \beta_1, \beta_1\}\!\}$
where the subscript denotes the vertex's degree. The two multisets
are different, so the two graphs are not $1$-WL equivalent.

This shows that $1$-WL sometimes distinguishes graphs with the same
vertex and edge counts. The pair $(2K_3, C_6)$ of Section~\ref{sec:warmup}
is the smallest non-trivial pair on which $1$-WL fails to
distinguish.
\end{example}

\subsection{Message-passing graph neural networks}\label{sec:mpnn}

A message-passing graph neural network takes a featured graph and
produces first a vertex representation and then a graph-level output.
The vertex representation at each layer is a function of the previous
layer's representation at the vertex and the multiset of previous-layer
representations at its neighbours. Aggregation over the neighbour
multiset uses a permutation-invariant function; the update at the
vertex is arbitrary. The graph-level output combines the final-layer
vertex representations into a single vector.

\begin{definition}[Message-passing GNN]\label{def:mpnn}
A message-passing graph neural network (MPNN) with $T$ layers is
determined by four functions: a feature embedding $\iota :
\mathcal{X} \to \R^{d_0}$, and for each layer $t \in \{1, \dots, T\}$
a permutation-invariant aggregation $\agg^{(t)}$ on multisets over
$\R^{d_{t-1}}$, an update $\upd^{(t)} : \R^{d_{t-1}} \times
\R^{d_{t-1}} \to \R^{d_t}$, and a permutation-invariant readout
$\readout$ on multisets over $\R^{d_T}$. Given a featured graph
$(\gph, x)$, the MPNN computes representations $h_v^{(t)} \in
\R^{d_t}$ layer by layer:
\[
h_v^{(0)} = \iota(x_v),
\quad
h_v^{(t)} = \upd^{(t)}\!\left(h_v^{(t-1)},\ \agg^{(t)}\!\left(\{\!\{h_u^{(t-1)} : u \in \neigh{v}\}\!\}\right)\right),
\]
for $t \in \{1, \dots, T\}$, and the graph-level output
\[
M(\gph, x) = \readout\!\left(\{\!\{h_v^{(T)} : v \in \vertices\}\!\}\right).
\]
\end{definition}

Write $\mpnn$ for the class of all such maps $M$ from featured graphs
to $\R$, over every choice of depth $T$, hidden widths $d_1, \dots,
d_T$, and continuous choice of $\iota$, $\agg^{(t)}$, $\upd^{(t)}$,
and $\readout$. Write $\mpnn_{\mathrm{sum}}$ for the subclass in
which every $\agg^{(t)}$ and $\readout$ is a sum aggregation composed
with a feed-forward network (a graph isomorphism network in the sense
of \citet{XuGIN2019}).

\begin{example}[A depth-$1$ MPNN on $C_6$]\label{ex:mpnn-c6}
Take $B = C_6$ with the alternating feature from Example~\ref{ex:wl-feat}.
Choose $\iota$ to be the identity (so features live in $\R$), sum
aggregation, an update $\upd^{(1)}(h_v, s) = h_v + s$, and a sum
readout. Then
\[
h_v^{(1)} = x_v + \sum_{u \in \neigh{v}} x_u.
\]
An odd vertex $v$ has $x_v = +1$ and two neighbours with features
$-1$, so $h_v^{(1)} = 1 + (-1) + (-1) = -1$. An even vertex has
$h_v^{(1)} = -1 + 1 + 1 = +1$. The readout returns
$M(B, x) = 3 \cdot (-1) + 3 \cdot (+1) = 0$.
\end{example}

\begin{example}[Depth-$2$ MPNN on the path $P_4$]\label{ex:mpnn-p4}
Take $P_4$ from Example~\ref{ex:wl-tree} with a constant feature $x_v = 1$
for every $v$. Use $\iota = \mathrm{id}$, sum aggregation at both
layers, updates $\upd^{(1)}(h, s) = h + s$ and $\upd^{(2)}(h, s) =
h + s$, and a sum readout. Compute layer by layer.

Round $0$: $h_v^{(0)} = 1$ for every $v$.

Round $1$: an endpoint ($1$ or $4$) has one neighbour, so its
representation is $h_1^{(1)} = h_4^{(1)} = 1 + 1 = 2$. An interior
vertex ($2$ or $3$) has two neighbours, so $h_2^{(1)} = h_3^{(1)} =
1 + 1 + 1 = 3$.

Round $2$: vertex $1$'s only neighbour is $2$, with $h_2^{(1)} = 3$,
so $h_1^{(2)} = 2 + 3 = 5$. Vertex $4$ mirrors: $h_4^{(2)} = 5$.
Vertex $2$'s neighbours are $1$ (with $h^{(1)} = 2$) and $3$ (with
$h^{(1)} = 3$), so $h_2^{(2)} = 3 + 2 + 3 = 8$. Vertex $3$ mirrors:
$h_3^{(2)} = 8$.

Readout: $M(P_4, 1) = 5 + 8 + 8 + 5 = 26$.
\end{example}

\begin{example}[Two graphs with the same MPNN output at
depth $1$ but different depth-$2$ outputs]\label{ex:mpnn-depth}
Consider two graphs on $6$ vertices with a constant feature $x_v =
1$. The first is $A = 2 C_3$, two disjoint triangles; the second is
$B = C_6$. Every vertex in either graph has degree $2$. With the
depth-$1$ MPNN from Example~\ref{ex:mpnn-c6} (sum aggregation, additive
update, sum readout), a vertex $v$ receives $h_v^{(1)} = 1 + 1 + 1 =
3$ on either graph, and $M(A, 1) = M(B, 1) = 6 \cdot 3 = 18$.

At depth $2$, a vertex $v$ receives $h_v^{(2)} = h_v^{(1)} +
\sum_{u \in \neigh{v}} h_u^{(1)} = 3 + 3 + 3 = 9$ on either graph
as well: every vertex's round-$1$ representation is $3$, and every
vertex has two neighbours. Readout: $M(A, 1) = M(B, 1) = 6 \cdot 9
= 54$. The two outputs still match.

This is consistent with $A \sim_{\wl} B$ (Section~\ref{sec:warmup}): $1$-WL
does not distinguish the two graphs, and the MPNN, whose output is
$1$-WL invariant by Proposition~\ref{prop:upper}, does not distinguish them
either.
\end{example}

An MPNN's output depends on the featured graph only through the
values passed by its aggregations and updates. If two featured graphs
are $1$-WL equivalent, the multiset of representations at each layer
agrees between them; the graph-level output agrees as well. This is
the content of the $1$-WL upper bound \citep{MorrisWL2019, XuGIN2019},
which Section~\ref{sec:dichotomy} states as one direction of the dichotomy.

\subsection{Two function classes}\label{sec:classes}

Fix a set $\mathcal{G}$ of featured graphs. Two subsets of real
functionals on $\mathcal{G}$ will be compared throughout the paper.

\begin{definition}[$1$-WL invariant functional]\label{def:wl-inv}
A real-valued map $\Pi : \mathcal{G} \to \R$ is \emph{$1$-WL
invariant} on $\mathcal{G}$ if $(\gph, x) \sim_{\wl} (\gph', x')$
implies $\Pi(\gph, x) = \Pi(\gph', x')$ for every $(\gph, x),
(\gph', x') \in \mathcal{G}$.
Write $\mathcal{F}_{\wl}(\mathcal{G})$ for the set of $1$-WL
invariant real functionals on $\mathcal{G}$.
\end{definition}

\begin{definition}[MPNN-computable functional]\label{def:mpnn-comp}
A real-valued map $\Pi : \mathcal{G} \to \R$ is
\emph{MPNN-computable} on $\mathcal{G}$ if for every $\varepsilon > 0$
there is an MPNN $M \in \mpnn_{\mathrm{sum}}$ with
$\sup_{(\gph, x) \in \mathcal{G}} |M(\gph, x) - \Pi(\gph, x)| <
\varepsilon$.
Write $\mathcal{F}_{\mpnn}(\mathcal{G})$ for the uniform closure on
$\mathcal{G}$ of the MPNN-computable functionals.
\end{definition}

$\mathcal{F}_{\mpnn}(\mathcal{G})$ is the set of functionals a
message-passing GNN can approximate arbitrarily well on
$\mathcal{G}$; $\mathcal{F}_{\wl}(\mathcal{G})$ is the set of
functionals that respect the equivalence relation $\sim_{\wl}$. The
dichotomy of Section~\ref{sec:dichotomy} states that these two sets
coincide, once one restricts to continuous functionals on a compact
featured-graph class of bounded size.
